% appendix-a.tex
\chapter{Repository und interaktive Dokumentation}
\label{chap:anhang_a}

\section{Aufbau des Repositories}

\begin{lstlisting}[caption={Verzeichnisstruktur des Projekt-Repositories.},label={lst:repo}]
carrera_rl/
|
+-- Notebook_prototyping.ipynb   zentrales interaktives Arbeitsbuch
+-- README.md                    Kurzeinstieg
+-- requirements.txt             Python-Abhaengigkeiten
|
+-- data/                        Streckenbild, Fahrzeug-Sprite, Splines (.npy),
|                                Demonstrationsdatensaetze (.npz)
+-- models/                      je Lauf ein Verzeichnis mit best_model.zip,
|                                train_config.json und evaluations.npz
+-- tensorboard_logs/            Diagnosegroessen je Lauf
+-- Quellen und Paper/           wissenschaftliche Literatur
+-- Other tools/
|   +-- Auto Parameter messen/   Videos und Auswertung der Fahrzeugmessung
|
+-- src/
    +-- envs/
    |   +-- carrera_2d_env.py       Gymnasium-Umgebung: Physik, Lidar, Kamera
    +-- utils/
    |   +-- tracklimits_cv.py       OpenCV-Extraktion von Raendern und Mittellinie
    |   +-- virtual_camera.py       globale und fahrzeugfeste Ansicht, Quantisierung
    |   +-- reward_func.py          Belohnungsfunktion (Version v3_rundenzeit)
    |   +-- framestapel_buffer.py   speicheroptimierter Replay-Buffer
    |   +-- framestapel_gespreizt.py  gespreizter Framestapel (Abstand 1, 2, 4, 8)
    |   +-- model_free.py           SensorSuite: Raycast-Berechnung
    |   +-- benchmark.py            Durchsatz- und Rechenaufwandsmessung
    +-- train/
    |   +-- lidar_experte.py        Training des privilegierten Lehrer-Agenten
    +-- eval/
        +-- vergleich.py            Nachauswertung aller Modelle, einheitlich
        +-- ablation_lidar.py       gezielte Verfaelschung der Beobachtung
        +-- bc_fahren.py            vortrainiertes Netz selbst fahren lassen
        +-- backbone_probe.py       Guete der uebertragenen Faltungsschichten
        +-- fahrlinien.py           Trajektorienvergleich Experte / Agent
        +-- rauchtest.py            schneller Funktionstest eines Modells
\end{lstlisting}

\section{Leitfaden durch das Notebook}

Das \texttt{Notebook\_prototyping.ipynb} umfasst \num{70} Zellen und bildet den Arbeitsverlauf in der Reihenfolge seiner tatsächlichen Entstehung ab. Tabelle~\ref{tab:notebook} ordnet die Zellbereiche den Abschnitten dieses Berichts zu.

\begin{table}[htbp]
  \centering
  \caption{Zuordnung der Notebook-Zellen zu den Abschnitten dieses Berichts.}
  \label{tab:notebook}
  \small
  \begin{tabularx}{\linewidth}{@{}l X l@{}}
    \toprule
    \textbf{Zellen} & \textbf{Inhalt} & \textbf{Abschnitt} \\
    \midrule
    0--2   & Zielsetzung und Brainstorming zu Eingabekonzepten & \ref{sec:ziele} \\
    3--5   & Streckenbild laden, Fahrzeug initialisieren, manuelles Fahren mit WASD & \ref{subsec:idee} \\
    6--11  & Streckengrenzen per OpenCV, interaktives HSV-Tuning, Mittelliniengenerierung & \ref{subsubsec:vorbereitung} \\
    12--15 & Zustände auf Basis der Grenzen, Start/Ziel; Belohnungs- und Raycast-Debugfenster & \ref{subsubsec:reward_entwurf} \\
    16--20 & Erstes RL-Training (PPO, MLP-Policy auf Raycasts), Agenten beobachten & \ref{subsubsec:erstes_training} \\
    21--26 & Virtuelle Kamera: fahrzeugfester Crop und globale Ansicht & \ref{subsubsec:vorbereitung} \\
    27--28 & Trainingszelle Raycasts & \ref{subsubsec:erstes_training} \\
    29--32 & Single-Frame, fahrzeugfest und global & \ref{subsubsec:eingabekonzepte} \\
    33--36 & Hybrid / Multi-Input (Bild und Skalarparameter) & \ref{subsubsec:eingabekonzepte} \\
    37--40 & Framestapel aus drei Bildern, fahrzeugfest und global & \ref{subsubsec:eingabekonzepte} \\
    41--43 & Beliebigen Agenten weitertrainieren & --- \\
    44--49 & Beobachtungszellen und Benchmark der Varianten & --- \\
    50--54 & Teacher-Student-Pipeline: Methodik, Datenaufzeichnung, CNN-Vortraining & \ref{sec:ergebnisse_pipeline} \\
    55--61 & SAC-Feintuning, große Trainingsläufe, PPO-Vergleichszelle & \ref{subsec:ppo_sac_vergleich} \\
    62--63 & Vergleich aller Trainingsläufe & \ref{sec:ergebnisse_quantitativ} \\
    64--65 & Agenten beobachten -- mit der Bildpipeline des Trainings & \ref{subsec:buffer} \\
    66--67 & Kontrollversuch: SAC auf Lidar unter \texttt{measured\_v2} & \ref{subsec:subpixel} \\
    68--69 & Beliebigen gespeicherten Agenten unter seiner Originalphysik beobachten & --- \\
    \bottomrule
  \end{tabularx}
\end{table}

\section{Wiederkehrende Aufrufe}

\begin{lstlisting}[language=bash,caption={Häufig benötigte Kommandos.},label={lst:kommandos}]
# virtuelle Umgebung aktivieren
source .rl_venv/bin/activate

# Trainingsverlaeufe ansehen
tensorboard --logdir ./tensorboard_logs/

# Lehrer-Agenten neu trainieren
python -m src.train.lidar_experte --steps 150000

# alle Modelle einheitlich nachauswerten
python -m src.eval.vergleich --episoden 10

# einzelnes Modell schnell pruefen
python -m src.eval.rauchtest --model models/<lauf>/best_model.zip

# Ablationsstudie auf der Lidar-Beobachtung
python -m src.eval.ablation_lidar
\end{lstlisting}

\noindent
Hinweis: Große Modellarchive unter \texttt{models/} können aufgrund ihrer Dateigröße nicht vollständig über Git verteilt werden.
